% ==========================================================================
% INFO 4617 Web Data Science -- Week 6 handout
% Strategy 6, cards that need to become rows: the Oscars loop from Chapter 6,
% built one step at a time.
%
% Every code block ran, in order, against the live page on 2026-09-24; each
% output block is the real output of the code above it. The two screenshots
% are real captures from the same day (see img/IMAGES.md).
%
% Build:  latexmk -pdf oscars-cards-to-rows.tex   (from this folder or the
%         repo root; the class is handouts/common/handout.cls)
% Check:  python check_outputs.py   (reruns every block against the live
%         page and compares each output with the one printed here)
% ==========================================================================
\makeatletter\def\input@path{{../common/}{handouts/common/}}\makeatother
\documentclass{handout}

\usepackage{accsupp}
\usetikzlibrary{decorations.pathreplacing}
\hypersetup{pdftitle={Cards to rows: the Oscars, one step at a time},
            pdfauthor={Brian C. Keegan}}

% a class name in running text: typewriter, may break after a hyphen
\DeclareUrlCommand\cls{\urlstyle{tt}}
% a numbered marker; the same number marks the same part of a card everywhere
\newcommand{\mk}[1]{\tikz[baseline=(m.base)]\node[circle,fill=tab-blue,
  text=white,inner sep=0pt,minimum size=1.2em,
  font=\sffamily\bfseries\scriptsize](m){#1};}
\tikzset{
  mkc/.style={circle,fill=tab-blue,text=white,inner sep=0pt,
              minimum size=11pt,font=\sffamily\bfseries\scriptsize},
  mkf/.style={mkc,minimum size=8.5pt,font=\sffamily\bfseries\tiny},
  thinlead/.style={tab-blue,line width=0.5pt,densely dotted},
  lead/.style={tab-blue,line width=0.8pt,densely dotted},
  brace/.style={tab-blue,line width=0.9pt,decorate,
                decoration={brace,amplitude=4pt}},
}

\title[INFO 4617 · Week 6 · Strategy 6 handout · checked against oscars.org on September 24, 2026]{Cards to rows: the Oscars, one step at a time}
\subtitle{Week 6 · Documents module · Textbook Chapter 6, Strategy 6}
\author{Brian C. Keegan, Ph.D.}
\institute{Department of Information Science, University of Colorado Boulder}
\date{}

\begin{document}

\maketitle

% ------------------------------------------------------------------ intro
\begin{frame}{What you will build}
\begin{columns}[T]
\begin{column}{0.605\textwidth}
The page for the 98th Academy Awards
(\url{https://www.oscars.org/oscars/ceremonies/2026}) lists 24 award
categories. It has no \texttt{<table>}. Each category is a \emph{card}: a
name, then one block for each honoree, first the winner and then the
nominees. Each honoree block names a person or group and a film
(Figure~1). This handout builds the loop from Wednesday's slide one piece
at a time and turns the cards into a DataFrame with one row per honoree:

\smallskip
{\fontsize{7.6}{9.2}\selectfont\setlength{\tabcolsep}{3.5pt}%
\begin{tabular}{@{}llll@{}}
\toprule
\hd{category} & \hd{honoree\_type} & \hd{entity} & \hd{film}\\
\midrule
Actor in a Leading Role & Winner & Michael B. Jordan & Sinners\\
Actor in a Leading Role & Nominees & Timoth\'ee Chalamet & Marty Supreme\\
\multicolumn{4}{@{}l}{\itshape \dots\ 125 rows in all, from 24 cards}\\
\bottomrule
\end{tabular}}

\smallskip
Each step adds one thing to the step before it. Run each step in your
notebook, and compare your output with the output here before you go to
the next step.

\smallskip
\begin{tabular}{@{}l@{\hspace{5pt}}l@{}}
{\color{cugold}\rule[-2pt]{2.4pt}{9pt}} & Code that you type.\\
{\color{handoutblockbody}\rule[-2pt]{2.4pt}{9pt}} & What Python shows back.\\
\mk{1}\,\dots\,\mk{7} & The same seven parts of a card, in\\
 & every figure and every step.
\end{tabular}
\end{column}
\begin{column}{0.375\textwidth}
\setlength{\unitlength}{\dimexpr\linewidth/1185\relax}% 1 px of the image; room for the markers at right
\BeginAccSupp{method=pdfstringdef,Alt={Screenshot of the Actor in a Leading Role card on oscars.org: the gold heading, then WINNER, MICHAEL B. JORDAN, Sinners, then NOMINEES and four more names with their films, then the next card, Actor in a Supporting Role. Numbered markers label the area around all the cards, one card, its name, one honoree, and the honoree's type, name, and film.}}%
\begin{tikzpicture}[x=\unitlength,y=-\unitlength]
  \node[anchor=north west,inner sep=0] at (0,0)
    {\includegraphics[width=1120\unitlength]{handouts/week-06/img/oscars_card.png}};
  \draw[black!25] (0,0) rectangle (1120,1190);
  \node[mkf] at (992,67) {3};            % the card's name
  \draw[thinlead] (156,172) -- (318,152); % Winner
  \node[mkf] at (345,150) {5};
  \node[mkf] at (360,212) {6};            % the person
  \draw[thinlead] (144,257) -- (318,273); % the film
  \node[mkf] at (345,275) {7};
  \draw[brace] (405,156) -- (405,280);    % one honoree
  \node[mkf] at (462,218) {4};
  \draw[brace] (1040,28) -- (1040,871);   % one card
  \node[mkf] at (1078,450) {2};
  % the area around all 24 cards; it continues below the picture
  \draw[tab-blue,line width=0.9pt] (1106,28) -- (1118,28) -- (1118,1150);
  \draw[tab-blue,line width=0.9pt,-latex] (1118,1150) -- (1118,1205);
  \node[mkf] at (1150,620) {1};
\end{tikzpicture}%
\EndAccSupp{}

\vspace{2pt}
{\small\raggedright \textbf{Figure 1.} One card as the browser shows it,
September 24, 2026. The numbers match Figure~2 in Step~3.\par}
\end{column}
\end{columns}
\end{frame}

% ------------------------------------------------------------------ step 1
\begin{frame}{Step 1 --- Ask for the page, and say who you are}
\begin{code}
import requests
from bs4 import BeautifulSoup
import pandas as pd

url = "https://www.oscars.org/oscars/ceremonies/2026"
headers = {"User-Agent": "Web Data Science/v1 brian.keegan@colorado.edu"}
response = requests.get(url, headers=headers)
response.status_code
\end{code}
\begin{out}
200
\end{out}
The \texttt{User-Agent} header tells the server who is asking: this course
and a contact address. Without it, oscars.org answers \texttt{403}
(Forbidden). \texttt{200} means that the server sent the page. In your own
projects, put your own name and address in the header.
\end{frame}

% ------------------------------------------------------------------ step 2
\begin{frame}{Step 2 --- Parse the HTML}
\begin{code}
soup = BeautifulSoup(response.text, "html.parser")
soup.title.get_text()
\end{code}
\begin{out}
'The 98th Academy Awards | 2026'
\end{out}
\texttt{soup} holds the page as a tree of tags that you can search. The
title tells you that you have the ceremony page, not an error page.
\end{frame}

% ------------------------------------------------------------------ step 3
\begin{frame}{Step 3 --- Find one card in DevTools}
Open the page in Chrome, right-click the name MICHAEL B. JORDAN, and
select \textbf{Inspect}. The Elements panel opens with the \texttt{<div>}
for that name selected (Figure~2). Each indented line sits inside the
nearest less-indented line above it: read up to find the containers around
the name, and down to find the parts next to it. Click a triangle to open
a line; Figure~2 has the film's \texttt{<div class="field">} open.

The \texttt{class} attribute on each line holds several class names,
separated by spaces; the line at \mk{5} has three: \cls{field},
\cls{field--name-field-honoree-type}, and \cls{winner}. oscars.org is
built with Drupal, a content management system that makes most of these
names from the site's own field names: \cls{field--name-} plus a field
name, and \cls{paragraph--type--} plus a type of block. These names tell
you what a div holds; names such as \cls{field} and \cls{field__item} tell
you only that the div is a box.

\begin{columns}[T]
\begin{column}{0.72\textwidth}
\setlength{\unitlength}{\dimexpr\linewidth/1040\relax}% 1 px of the image
\BeginAccSupp{method=pdfstringdef,Alt={Screenshot of the Chrome DevTools Elements panel for the Actor in a Leading Role card. Nested div tags, each with a long class attribute: the categories area, one category card with the name Actor in a Leading Role, the honorees list, and the first honoree with Winner, Michael B. Jordan (selected), and Sinners. Four more honorees are collapsed below. Numbered markers 1 to 7 point to the seven class names that the code uses.}}%
\begin{tikzpicture}[x=\unitlength,y=-\unitlength]
  \node[anchor=north west,inner sep=0] at (0,0)
    {\includegraphics[width=1040\unitlength]{handouts/week-06/img/oscars_devtools.png}};
  \draw[black!25] (0,0) rectangle (1040,824);
  \foreach \n/\y/\x in {1/14/90, 2/96/132, 3/151/160, 4/287/195,
                        5/341/222, 6/397/237, 7/562/243} {
    \draw[lead] (46,\y) -- (\x,\y);
    \node[mkc] at (30,\y) {\n};
  }
  % what you right-clicked
  \node[overlay,fill=white,rounded corners=2pt,inner sep=2pt,text=tab-blue,
        font=\sffamily\bfseries\scriptsize,anchor=west] at (905,451)
        {$\leftarrow$ you clicked};
  % the four collapsed honorees
  \draw[brace] (596,688) -- (596,797);
  \node[anchor=west,text=tab-blue,font=\sffamily\bfseries\scriptsize]
    at (612,742) {4 more honorees \tikz[baseline=-0.6ex]\node[mkc]{4};};
\end{tikzpicture}%
\EndAccSupp{}
\end{column}
\begin{column}{0.26\textwidth}
\small\raggedright
\textbf{What each part holds}\par\smallskip
\begin{tabular}{@{}l@{\hspace{4pt}}>{\raggedright\arraybackslash}p{0.82\linewidth}@{}}
\mk{1} & the area that holds all 24 cards\\[2pt]
\mk{2} & one card (a category)\\[2pt]
\mk{3} & the card's name\\[2pt]
\mk{4} & one honoree, a winner or a nominee\\[2pt]
\mk{5} & \texttt{Winner} or \texttt{Nominees}\\[2pt]
\mk{6} & the person or group (the \emph{entity})\\[2pt]
\mk{7} & the film\\
\end{tabular}

\vspace{8pt}
\textbf{Figure 2.} Chrome DevTools, Elements panel, after
\textbf{Inspect} on MICHAEL B. JORDAN. September 24, 2026.
\end{column}
\end{columns}
\end{frame}

% ------------------------------------------------------------------ step 4
\begin{frame}{Step 4 --- Match one class name, then count the matches}
\begin{code}
name_tag = soup.find("div", class_="field--name-field-award-category-oscars")
name_tag["class"]
\end{code}
\begin{out}
['field',
 'field--name-field-award-category-oscars',
 'field--type-entity-reference',
 'field--label-hidden',
 'field__item']
\end{out}
BeautifulSoup keeps the \texttt{class} attribute as a list of names.
\texttt{class\_="..."} matches a tag when your name is any one item in
that list, so you give only the one name that you need. (The underscore
is necessary because \texttt{class} is a reserved word in Python.) Now
count what each name matches on the whole page. The last item is the
whole attribute from the line at \mk{5}, copied from DevTools:
\begin{code}
for name in ["field", "field__item", "paragraph",
             "field--name-field-award-categories", "paragraph--type--award-category",
             "field--name-field-award-category-oscars", "paragraph--type--award-honoree",
             "field--name-field-honoree-type", "field--name-field-award-entities",
             "field--name-field-award-film",
             "field field--name-field-honoree-type winner"]:
    print(len(soup.find_all("div", class_=name)), name)
\end{code}
\begin{dfout}{@{}l@{\hspace{14pt}}c@{}}
\ol|693 field| & \\
\ol|445 field__item| & \\
\ol|154 paragraph| & \\
\ol|1 field--name-field-award-categories| & \mk{1}\\
\ol|24 paragraph--type--award-category| & \mk{2}\\
\ol|24 field--name-field-award-category-oscars| & \mk{3}\\
\ol|125 paragraph--type--award-honoree| & \mk{4}\\
\ol|125 field--name-field-honoree-type| & \mk{5}\\
\ol|125 field--name-field-award-entities| & \mk{6}\\
\ol|125 field--name-field-award-film| & \mk{7}\\
\ol|25 field field--name-field-honoree-type winner| & \\
\end{dfout}
The general names match hundreds of divs across the page; the specific
names match only what you want. Check each count against the page: it
shows 24 category headings, so 24 cards \mk{2} is correct, and 125
honorees \mk{4} with 125 of each field \mk{5}\,\mk{6}\,\mk{7} means one of
each per honoree. The whole attribute finds only the 25 winners, because a
string with spaces must match the attribute exactly, names and order.
Select the one name that describes the content, then check its count.
\end{frame}

% ------------------------------------------------------------------ step 5
\begin{frame}{Step 5 --- Go to the area \mk{1}, then collect the cards \mk{2}}
\begin{code}
area = soup.find("div", class_="field--name-field-award-categories")
categories = area.find_all("div", class_="paragraph--type--award-category")
len(categories)
\end{code}
\begin{out}
24
\end{out}
\texttt{find} returns the first match: one tag. \texttt{find\_all} returns
a list of all the matches. You call \texttt{find\_all} on \texttt{area},
not on \texttt{soup}, so it searches only inside the awards area.
\texttt{categories} is a list of the 24 cards, in the order of the page.
\end{frame}

% ------------------------------------------------------------------ step 6
\begin{frame}{Step 6 --- Take one card and read its name \mk{3}}
\begin{code}
cat = categories[0]
name_tag = cat.find(class_="field--name-field-award-category-oscars")
name_tag.get_text(strip=True)
\end{code}
\begin{out}
'Actor in a Leading Role'
\end{out}
Build the steps for one card before you build them for 24.
\texttt{cat.find(...)} searches only inside this card. This \texttt{find}
has no \texttt{"div"}, so it matches any tag with the class; every field
on this page is a div, so both forms work. \texttt{get\_text(strip=True)}
gives the text inside the tag, without the spaces and line breaks around
it.
\end{frame}

% ------------------------------------------------------------------ step 7
\begin{frame}{Step 7 --- Collect the honorees \mk{4} in that one card}
\begin{code}
honorees = cat.find_all("div", class_="paragraph--type--award-honoree")
len(honorees)
\end{code}
\begin{out}
5
\end{out}
In Step~4, this class name matched 125 divs on the whole page. Searched
from \texttt{cat}, it matches the 5 in this card. Where you start a search
decides what you get back. \texttt{find\_all} also looks at every level
below \texttt{cat}, so you do not step through the two wrapper divs between
the card and its honorees (\cls{field--name-field-award-honorees} and
\cls{field__item} in Figure~2).
\end{frame}

% ------------------------------------------------------------------ step 8
\begin{frame}{Step 8 --- Extract the three fields \mk{5}\,\mk{6}\,\mk{7} from one honoree}
\begin{code}
h = honorees[0]
type_tag = h.find(class_="field--name-field-honoree-type")
entity_tag = h.find(class_="field--name-field-award-entities")
film_tag = h.find(class_="field--name-field-award-film")
print(type_tag.get_text(strip=True))
print(entity_tag.get_text(strip=True))
print(film_tag.get_text(strip=True))
\end{code}
\begin{out}
Winner
Michael B. Jordan
Sinners
\end{out}
Each \texttt{find} starts from \texttt{h}, so it finds only the fields of
this honoree. To see what \texttt{strip=True} removes, look at the raw
text:
\begin{code}
entity_tag.text
\end{code}
\begin{out}
'\nMichael B. Jordan\n'
\end{out}
The line breaks come from the HTML around the inner \cls{field__item} div.
The type is also a class name (\cls{winner} or \cls{nominee}), but the
text is what you want in the table.
\end{frame}

% ------------------------------------------------------------------ step 9
\begin{frame}{Step 9 --- Make one row for each honoree in the card}
\begin{code}
category_name = name_tag.get_text(strip=True)
rows = []
for h in honorees:
    rows.append({
        "category": category_name,
        "honoree_type": h.find(class_="field--name-field-honoree-type").get_text(strip=True),
        "entity": h.find(class_="field--name-field-award-entities").get_text(strip=True),
        "film": h.find(class_="field--name-field-award-film").get_text(strip=True),
    })
rows[1]
\end{code}
\begin{out}
{'category': 'Actor in a Leading Role',
 'honoree_type': 'Nominees',
 'entity': 'Timothée Chalamet',
 'film': 'Marty Supreme'}
\end{out}
The three \texttt{find} calls from Step~8 go inside the loop, and their
text goes into a dictionary: one dictionary per row, keys as column names
(\texttt{rows[1]} is the second of 5). \texttt{category\_name}, from
Step~6, puts the card's name in each row. The page shows NOMINEES once,
but the HTML has the label in every nominee block: a script hides the
repeats in the browser (\texttt{style="display: none;"} in DevTools), and
\texttt{requests} does not run scripts, so every row gets its label.
\end{frame}

% ----------------------------------------------------------------- step 10
\begin{frame}{Step 10 --- Put the card loop inside a loop over all 24 cards}
\begin{code}
rows = []
for cat in categories:                                                      # Step 5
    name_tag = cat.find(class_="field--name-field-award-category-oscars")   # Step 6
    if name_tag is None:                                                    # guard
        continue
    category_name = name_tag.get_text(strip=True)                           # Step 6
    for h in cat.find_all("div", class_="paragraph--type--award-honoree"):  # Step 7
        type_tag = h.find(class_="field--name-field-honoree-type")          # Step 8
        entity_tag = h.find(class_="field--name-field-award-entities")      # Step 8
        film_tag = h.find(class_="field--name-field-award-film")            # Step 8
        rows.append({                                                       # Step 9
            "category": category_name,
            "honoree_type": type_tag.get_text(strip=True) if type_tag else None,
            "entity": entity_tag.get_text(strip=True) if entity_tag else None,
            "film": film_tag.get_text(strip=True) if film_tag else None,
        })
len(rows)
\end{code}
\begin{out}
125
\end{out}
Two things are new. The outer loop \texttt{for cat in categories:} runs
Steps~6 to~9 once per card, and \texttt{rows = []} moves above it so all
24 cards add to one list. And guards: \texttt{find} returns \texttt{None}
when nothing matches, and \texttt{.get\_text()} on \texttt{None} raises an
\texttt{AttributeError} that stops the loop. \texttt{if name\_tag is None:
continue} skips a card with no name; \texttt{... if type\_tag else None}
records a missing field as \texttt{None}. No guard fires on the 2026 page;
the guards keep the loop running if one card is shaped differently.
\end{frame}

% ----------------------------------------------------------------- step 11
\begin{frame}{Step 11 --- Make the DataFrame, then check it against the page}
\begin{code}
df = pd.DataFrame(rows)
df.shape
\end{code}
\begin{out}
(125, 4)
\end{out}
\begin{code}
df.head()
\end{code}
\begin{dfout}{@{}rllll@{}}
 & \hd{category} & \hd{honoree\_type} & \hd{entity} & \hd{film}\\
\midrule
\textbf{0} & Actor in a Leading Role & Winner & Michael B. Jordan & Sinners\\
\textbf{1} & Actor in a Leading Role & Nominees & Timoth\'ee Chalamet & Marty Supreme\\
\textbf{2} & Actor in a Leading Role & Nominees & Leonardo DiCaprio & One Battle after Another\\
\textbf{3} & Actor in a Leading Role & Nominees & Ethan Hawke & Blue Moon\\
\textbf{4} & Actor in a Leading Role & Nominees & Wagner Moura & The Secret Agent\\
\end{dfout}
\begin{code}
df["honoree_type"].value_counts()
\end{code}
\begin{out}
honoree_type
Nominees    100
Winner       25
Name: count, dtype: int64
\end{out}
25 winners for 24 categories looks like a bug in the loop. Look at the
winners of one category:
\begin{code}
df[(df["category"] == "Live Action Short Film") & (df["honoree_type"] == "Winner")]
\end{code}
\begin{dfout}{@{}rllll@{}}
 & \hd{category} & \hd{honoree\_type} & \hd{entity} & \hd{film}\\
\midrule
\textbf{70} & Live Action Short Film & Winner & Sam A. Davis and Jack Piatt & The Singers\\
\textbf{71} & Live Action Short Film & Winner & Alexandre Singh and Natalie Musteata & Two People Exchanging Saliva\\
\end{dfout}
The category ended in a tie (the first Oscars tie since 2013), so the
table is correct: check a surprising number against the page before you
change your code. Check what each column holds, too. You named the
columns, but the site decides what goes in each field: for International
Feature Film, \texttt{entity} is the film and \texttt{film} is the country
(\emph{Sentimental Value}, Norway); for Music (Original Song),
\texttt{film} is the song title.

\begin{alertblock}{The whole path}
\texttt{soup} $\rightarrow$ area \mk{1} $\rightarrow$ each card \mk{2}
$\rightarrow$ its name \mk{3} and honorees \mk{4} $\rightarrow$ each
honoree's type \mk{5}, entity \mk{6}, and film \mk{7} $\rightarrow$ one
dictionary per row $\rightarrow$ \texttt{pd.DataFrame(rows)}. If your
counts differ, the page has changed: inspect a card again (Figure~2).
\end{alertblock}
\end{frame}

\end{document}
